\section{Szczegółowa specyfikacja projektu i opis wymagań}

\subsection{Cel i kryterium sukcesu}

Sukces projektu definiowany jest jako uzyskanie klasyfikatora (lub metody statystycznej) osiągającego \textbf{dokładność co najmniej $51\%$} na wcześniej odłożonym zbiorze testowym, wraz z opisem warunków, w których wynik jest miarodajny (długość tekstu, język, źródło generatora).

\subsection{Wejście / wyjście}

\begin{itemize}
	\item \textbf{Wejście:} tekst w formie ciągłej (plik lub łańcuch znaków) spełniający minimalną długość.
	\item \textbf{Wyjście:} etykieta \{\texttt{human}, \texttt{ai}\} oraz -- jeśli to możliwe -- wynik liczbowy (score / prawdopodobieństwo).
\end{itemize}

\subsection{Założenia}

\begin{itemize}
	\item Prace rozpoczynają się od podejścia statystycznego (zgodnie z konsultacją eksperta przy małej skali danych).
	\item Dla wariantu ML teksty treningowe powinny być możliwie zbliżone tematycznie (rekomendacja eksperta).
	\item Limit długości tekstu zostanie ustalony empirycznie.
\end{itemize}

\subsection{Ograniczenia}

Dokładność $51\%$ jest progiem minimalnym akceptacji, nie celem optymalnym. Wynik może nie przenosić się na inne domeny, modele LLM ani język polski, jeśli trening odbywał się głównie na danych angielskich.
